% Responsable: cada quien sus cifras (bloques separados para no chocar).
% REGLA: toda cifra del texto sale de aquí y cada macro dice de qué CSV viene.
% Si se repite una corrida, se cambia SOLO aquí y el paper queda coherente.
% Formato: coma decimal, punto de miles (como en las diapositivas).

% ===== Línea base (Martín) =============================================
% Fuente: Modelos/04_linea_base.py
\newcommand{\cMcero}{0,0387}          % M0 inicial (exacto: 0,038684)
\newcommand{\cHogares}{3.732}         % hogares de Antioquia en la muestra
\newcommand{\cIndicadores}{15}
\newcommand{\cMceroExacto}{0,038684}
\newcommand{\cPobresHog}{523}           % hogares pobres en la línea base
\newcommand{\cHcero}{9,45\,\%}         % verificado en CLAUDE.md de Modelo3 §3
\newcommand{\cAcero}{40,94\,\%}

% ===== Modelos 1 y 2 (Martín) ==========================================
% Fuente: \pendiente{CSV de resultados M1/M2 en el escenario común}

% ===== Modelos 3 y 4 (Jerónimo) ========================================
% Tamaño del modelo en la VERSIÓN INICIAL (costos por persona, meta 10 %, réplica 0). Fuente: docs_proyecto/revision-formulacion-modelo3.md
% y CLAUDE.md de Modelo3 §3 y §5 (meta 10 %, réplica 0)
\newcommand{\cParesActivos}{8.602}     % pares (i,j) con M_ij = 1
\newcommand{\cParesMtres}{45}          % pares (indicador, zona), versión inicial (con rezago)
\newcommand{\cParesMcuatro}{53}        % pares (indicador, celda) en M4, escenario final (m34_tamano_modelo.py)
\newcommand{\cCortesMtres}{630}
\newcommand{\cCortesMcuatro}{324 y 524}   % escenario final, 100 réplicas
\newcommand{\cBinPaper}{171.672}      % Lin 6-10 literal, b1 para todos los hogares
\newcommand{\cBinDirecta}{9.125}       % N = b2 y b1 solo para pobres: 8.602 + 523
\newcommand{\cRestrDirecta}{27.376}
\newcommand{\cBinCortes}{1.153}
\newcommand{\cRestrCortes}{1.615}
\newcommand{\cTiempoCortes}{2,1~s}
% Escenario final (datos-finales-10nov), 100 réplicas. Fuente: m34_tamano_modelo.py
\newcommand{\cParesMtresFin}{42}
\newcommand{\cCortesMtresFin}{318 y 542}
\newcommand{\cBinMtresFin}{841 a 1.065}
\newcommand{\cTiempoMtresMed}{1~s}
\newcommand{\cTiempoMtresMax}{22~s}
\newcommand{\cLogroFijos}{1.647}       % pares de logro que el M1 no puede quitar en un año
% Modelo 4. Fuente: CLAUDE.md de Modelo3 §8 y §7.4 (clústeres no dependen de los costos)
\newcommand{\cDiscAcuZona}{14,5}
\newcommand{\cDiscAcuClu}{4,1}
\newcommand{\cSeparaPobreza}{4,3}
\newcommand{\cSiluetaMin}{0,26}
\newcommand{\cSiluetaMax}{0,28}
\newcommand{\cDiscInasClu}{51,7}
\newcommand{\cHcabecera}{6,2\,\%}          % H0 de la cabecera, CLAUDE.md de Modelo3 §3
% Validación del escenario común contra el M1/M2 de main (10-nov). Fuente: resultados/m34/validacion_10nov.csv
\newcommand{\cValidOK}{18 pruebas}
\newcommand{\cParesPobres}{2.592}
\newcommand{\cVcincoVeinte}{5,9\,\%}
\newcommand{\cVcincoTreinta}{2,7\,\%}
% Validación 12-nov (granularidad máxima con meta por zona y r_B). Fuente: resultados/m34/validacion_12nov.csv
\newcommand{\cGunoTreinta}{0,4\,\%}
\newcommand{\cGdosMin}{0,3}
\newcommand{\cGdosMax}{1,0\,\%}
% Corridas del 10-nov (escenario común contra el M1/M2 de main), CONGELADAS. Fuente: resultados/m34/*_10nov.csv
% en la rama Modelo3, etiqueta datos-finales-10nov. Costos: 100 réplicas; equidad: 100 (B <= 10) y 30 (B >= 25).
\newcommand{\cReplicas}{30}               % equidad con B >= 25
\newcommand{\cReplicasCosto}{100}
\newcommand{\cCorridasCosto}{4.500}
\newcommand{\cMtresIPmin}{6}
\newcommand{\cMtresIPmax}{17}
\newcommand{\cMtresIPalta}{3,5 a 3,9}
\newcommand{\cMcuatroPlacMin}{14\,\%}
\newcommand{\cMcuatroPlacMax}{31\,\%}
\newcommand{\cReplicasBajo}{100}
\newcommand{\cMtresPctMin}{13\,\%}
\newcommand{\cMtresPctMax}{65\,\%}
\newcommand{\cMcuatroPlacEq}{9--10\,\%}
\newcommand{\cInfactiblesRho}{900}
% ===== Comparación principal (12-nov): restricciones de equidad de Martín tal cual =====
% Criterio de costo con meta de M0 en cada zona. Fuente: resultados/m34/costo_meta_zona_crudo_12nov.csv e
% ic_costo_meta_zona_12nov.csv (m34_tablas_12nov.py); IP = info_perfecta_meta_zona_12nov.csv (M2 sin E-7).
\newcommand{\zReplicas}{54}
\newcommand{\zCorridas}{2.430}
\newcommand{\zGapMax}{$8{,}8\times10^{-7}$}
\newcommand{\zMtresIPmin}{3,2}            % M3/IP, mínimo (meta 28 %)
\newcommand{\zMtresIPmax}{12,5}           % M3/IP, máximo (meta 2 %)
\newcommand{\zMtresIPbaja}{5,5 a 12,5}    % metas del 2 al 10 %
\newcommand{\zMtresIPalta}{3,2 a 3,9}     % metas del 20 al 30 %
\newcommand{\zMcuatroIPmin}{3,4}
\newcommand{\zMcuatroIPmax}{8,3}
\newcommand{\zMcuatroMtresMenor}{5}       % metas con IC de M4/M3 < 1 (2, 12, 20, 22, 24 %)
\newcommand{\zMcuatroMtresMayor}{2}       % metas con IC de M4/M3 > 1 (28, 30 %)
\newcommand{\zMcuatroMtresIgual}{8}
\newcommand{\zPlacMin}{14\,\%}            % M4 más barato que el placebo, desde la meta del 6 %
\newcommand{\zPlacMax}{37\,\%}
\newcommand{\zInclMtres}{27 y 55\,\%}     % inclusión mediana por meta, M3
\newcommand{\zInclMcuatro}{21 y 49\,\%}
\newcommand{\zExclMtres}{56 y 93\,\%}
\newcommand{\zIPzonaDiez}{5,17}           % IP (M2) con meta del 10 % por zona
\newcommand{\cIPglobalDiez}{1,32}         % IP (M2) con meta global del 10 % (10-nov)
% Criterio de cobertura con r_B. Fuente: resultados/m34/cobertura_rB_crudo_12nov.csv y rmax_12nov.csv
% (m34_rmax_12nov.py; ic_rmax_12nov.csv con m34_tablas_12nov.py), 100 réplicas.
\newcommand{\zCorridasCob}{1.350}           % 1.170 con r_B y 180 con r_B y rho_B (B = 25 y 50)
\newcommand{\zCorridasRmax}{1.800}
\newcommand{\zRmaxFracMin}{0}             % r_max(M3)/r_max(IP), mediana, B = 2
\newcommand{\zRmaxFracMax}{56\,\%}        % B = 100
